\section{Related Work}
\label{sec:related}

\paragraph{Classical plant segmentation.} Colour indices such as excess green, together with
thresholds in HSV or CIELAB space, are the standard way to separate vegetation from soil in field
images~\cite{woebbecke1995color,hamuda2016survey}. They are fast and easy to interpret but depend
on lighting and soil colour, and they cannot distinguish lesion types that have similar colours,
such as brown pitting and soil seen through a hole.

\paragraph{Deep learning for plant phenotyping.} Pretrained CNNs and vision transformers are widely
used to classify or regress plant traits from images~\cite{kamilaris2018deep}. Self-supervised
backbones such as DINOv2~\cite{oquab2023dinov2} and DINOv3~\cite{dinov3} provide strong frozen
features, so only a small head has to be trained, which suits small labelled sets. These models
are usually applied at around $224\times224$ pixels, however, which is far too low for
millimetre-sized lesions on seedlings in a full field image.

\paragraph{Multiple instance learning.} In MIL, a label is given for a bag of instances but not
for the instances themselves~\cite{dietterich1997solving,carbonneau2018multiple}. Attention-based
MIL (ABMIL)~\cite{ilse2018attention} learns the pooling weights from the instance features and is
the usual deep MIL baseline, for example for gigapixel pathology slides, which share our problem of
small relevant regions in very large images. Our setting is a regression variant: the bag label is
an average damage over plants, which suggests a fixed weighted mean instead of learned attention.

\paragraph{Learning from comparisons.} When absolute labels are noisy, relative judgements such as
``image A is more damaged than image B'' can be more reliable. Pairwise ranking losses, going back
to RankNet~\cite{burges2005learning}, train a model to order pairs correctly and can be combined
with a regression loss.

\paragraph{Domain shift in agricultural vision.} Models trained on one field, season or set of
lighting conditions often degrade sharply elsewhere. For example, a plant disease classifier with
over 99\,\% accuracy on laboratory images dropped to about 31\,\% on field images~\cite{mohanty2016using},
and multi-site datasets such as the Global Wheat Head Dataset were created to measure
this effect~\cite{david2020global}. We therefore evaluate on unseen trials without any adaptation.
